\section{Canonical analytic division retains its exact formal quotients and filtration}

\begin{theorem}[A whole owner-supported relation is zero componentwise]
\label{hir:actual-owner-supported-formal-relation-uniqueness}
\uses{hir:actual-formal-initial-exponent-coefficient-product,hir:actual-monic-vertex-family-whole-division}
\lean{Hironaka.HironakaActualOwnerAnalyticDivision.owner_supported_relation_eq_zero}
\leanok
For any genuine monic vertex family of a whole formal ideal, a whole
relation $\sum_a G_aH_a=0$ whose $a$th quotient is supported in its actual
disjoint monomial-owner part has every $H_a=0$. No convergence or inverse
operator hypothesis is needed.
\end{theorem}
\begin{proof}
Suppose some quotient is nonzero. Among the finitely many nonzero
quotients choose $a$ for which the shifted genuine first exponent
$c=a+\operatorname{inexp}(H_a)$ is least in the same monomial order.
The genuine first exponent of $G_aH_a$ is $c$, and its coefficient is
the nonzero first coefficient of $H_a$, since $G_a$ is monic at $a$.
Every other nonzero product has first exponent at least $c$. Equality
would give two quotient first coefficients in the same actual owner
part, forcing their vertex labels equal. Thus every different product
has strictly larger first exponent and its coefficient at $c$ is zero.
The whole sum therefore has the nonzero coefficient of $G_aH_a$ at $c$,
contrary to the relation. Empty families are included.
\end{proof}

\begin{theorem}[Every actual owner division has the literal recursively constructed formal quotients]
\label{hir:actual-owner-division-formal-quotient-dictionary}
\uses{hir:actual-owner-supported-formal-relation-uniqueness,hir:actual-monic-vertex-family-whole-division,hir:actual-monic-family-canonical-remainder}
\lean{Hironaka.HironakaActualOwnerAnalyticDivision.owner_division_quotients_eq_familyFormalQuotient}
\leanok
A whole owner-supported division with the genuine canonical normal
remainder has precisely the previously constructed recursive formal
quotients, for the same monic vertex family and the same numerator.
\end{theorem}
\begin{proof}
The genuine recursion gives another whole division with owner-supported
quotients. Its remainder equals the canonical normal remainder by the
whole remainder uniqueness theorem. Subtract the two actual whole
identities. The resulting difference quotients are owner-supported:
outside their owner's part both original coefficients are zero. Apply
the preceding relation theorem to their whole relation, giving zero
in every difference component, exactly equality of every quotient.
\end{proof}

\begin{theorem}[Original canonical analytic division has the exact Taylor quotients and shifted order bounds]
\label{hir:actual-original-canonical-analytic-filtered-division}
\uses{hir:actual-owner-division-formal-quotient-dictionary,hir:actual-original-canonical-reduced-analytic-division,hir:actual-monic-family-centre-power-division,hir:actual-original-formal-taylor-maximal-power-kernel}
\lean{Hironaka.HironakaActualOwnerAnalyticDivision.exists_original_filtered_reduced_vertex_division}
\leanok
Every original ideal and original numerator internally give analytic
quotients by its same canonical analytic reduced family and the same
canonical analytic normal remainder. Every quotient's original Taylor
series is literally the genuine formal recursive quotient. If the
numerator belongs to the whole maximal-ideal power $\mathfrak m^q$,
the quotient at vertex $a$ belongs to
$\mathfrak m^{\max(q-|a|,0)}$.
\end{theorem}
\begin{proof}
Use the original analytic owner division, which internally proves
convergence of its actual quotients and remainder. Apply original Taylor
to its whole identity; the canonical analytic generators have exactly
the same canonical reduced formal series and its remainder has exactly
the canonical formal normal representative. Its supplied genuine owner
support and the preceding theorem therefore identify every Taylor
quotient with the literal formal recursion. Degree-first monicity of
each canonical reduced member implies all its coefficients below total
degree $|a|$ vanish. The genuine recursive shifted-weight division bound,
with every coordinate weight one, now says every coefficient of the
$a$th quotient below degree $q-|a|$ is zero. The actual original
Taylor/whole maximal-power dictionary converts this coefficient
vanishing into the stated whole analytic ideal power. It includes
$q=0$ and negative differences by the ordinary zero power convention.
\end{proof}

\begin{theorem}[Approximate original canonical relations lift to genuine analytic syzygies]
\label{hir:actual-original-canonical-analytic-syzygy-lifting}
\uses{hir:actual-original-canonical-analytic-filtered-division}
\lean{Hironaka.HironakaActualOwnerAnalyticDivision.exists_original_canonical_syzygy_of_relation_error}
\leanok
For an original ideal with its same canonical analytic reduced family,
any actual analytic coefficient family $P_a$ whose actual relation error
$E=\sum_aG_aP_a$ lies in $\mathfrak m^q$ has a genuine analytic syzygy
$C_a$ with
\[
 \sum_aG_aC_a=0,\qquad
 P_a-C_a\in\mathfrak m^{\max(q-|a|,0)}\quad\hbox{for every }a.
\]
The input is the actual approximate relation at a stated finite order;
no exact syzygy, relation lift, convergence or prepared tracking
certificate is supplied.
\end{theorem}
\begin{proof}
The actual whole error belongs to the original ideal, since every $G_a$
does. Divide it by that same canonical analytic family using the preceding
theorem. Its canonical normal remainder is zero: the entire original
normal projection has kernel exactly the original ideal. Thus
$E=\sum_aG_aQ_a$ with every actual analytic quotient $Q_a$ in
$\mathfrak m^{\max(q-|a|,0)}$. Set $C_a=P_a-Q_a$. Subtract the two whole
expressions for $E$; distributivity over the finite actual vertex set
proves $\sum_aG_aC_a=0$. Since $P_a-C_a=Q_a$, all the prescribed whole
power bounds follow. In particular this corrects any actual polynomial
relation among homogeneous initial forms whose original relation error
has the corresponding next-order vanishing.
\end{proof}

These are exact canonical analytic relation lifts. Turning them into
exceptional torsion-freeness of the normalized transformed generator ideal,
and proving a reverse current Hilbert--Samuel persistence criterion,
remain separate assertions.

\begin{theorem}[Actual homogeneous initial relations have prescribed analytic syzygy lifts]
\label{hir:actual-original-homogeneous-initial-analytic-syzygy}
\uses{hir:actual-original-canonical-analytic-syzygy-lifting,hir:actual-original-formal-taylor-maximal-power-kernel}
\lean{Hironaka.HironakaActualOwnerAnalyticDivision.exists_original_initial_syzygy_lift}
\leanok
Fix an original ideal with its actual canonical analytic generators $G_a$,
and an integer $D$ at least every vertex degree $|a|$. Let $P_a$ be actual
original analytic germs whose entire Taylor series are homogeneous of
degree $D-|a|$, and suppose the literal homogeneous initial relation
\[
 \sum_a T(P_a)\,\operatorname{hom}_{|a|}T(G_a)=0
\]
holds. Then internally produce actual analytic germs $C_a$ with
$\sum_aG_aC_a=0$ and
\[
 \operatorname{hom}_{D-|a|}T(C_a)=T(P_a),\qquad
 P_a-C_a\in\mathfrak m^{D+1-|a|}
\]
for every vertex. The input is a genuine homogeneous relation on the
same source initial forms; no vanishing order of its entire original
relation error or analytic syzygy is supplied.
\end{theorem}
\begin{proof}
Put $E=\sum_aG_aP_a$. For every exponent $c$ of degree at most $D$,
expand each product coefficient in its genuine finite antidiagonal.
A nonzero coefficient of $T(P_a)$ at $u$ forces $|u|=D-|a|$.
Consequently the complementary exponent $v$, with $u+v=c$, has
$|v|\le|a|$. If this inequality is strict, degree-first monicity of
the original canonical reduced generator makes its coefficient zero;
if equality holds, its coefficient is precisely that of its degree
$|a|$ homogeneous component. Thus every coefficient of $T(E)$ through
degree $D$ equals the same coefficient of the displayed whole zero
relation. All such coefficients vanish. The actual original maximal-power
dictionary therefore gives $E\in\mathfrak m^{D+1}$ internally.
Apply the preceding actual analytic syzygy correction with $q=D+1$.
It supplies $P_a-C_a\in\mathfrak m^{D+1-|a|}$ and the genuine whole
syzygy. Its Taylor correction has zero coefficients in degree $D-|a|$,
since $|a|\le D$. Taking that actual homogeneous component leaves
$T(P_a)$ unchanged; the latter series is homogeneous of exactly that
degree by hypothesis. This proves both the prescribed initial component
and every whole analytic correction power.
\end{proof}
